%MB9T5
\begin{dang}{HIỆN TƯỢNG VÀ ĐIỀU KIỆN XẢY RA SÓNG DỪNG}
% \textbf{Tìm số nút - số bụng giữa hai điểm bất kì}
%\begin{bang}[tabularx={X|X|X|X|X|X}]{}
%Hai đầu&A nút\newline B nút&A bụng \newline B bụng&A nút \newline B bụng&A nút \newline B chưa biết&A bụng \newline B chưa biết\\\hline
%\[\text{Số nút}\]&\[Sb+1\]&\[\dfrac{AB}{0,5\lambda}\]&\[\dfrac{AB}{0,5\lambda}+0,5\]&\[\text{Tách}\]&\[\text{Tách}\]\\\hline
%\[\text{Số bụng}\]&\[\dfrac{AB}{0,5\lambda}\]&\[Sn+1\]&\[\dfrac{AB}{0,5\lambda}+0,5\]&\[\text{Tách}\]&\[\text{Tách}\]\\
%\end{bang}
%\newghichu{
%\begin{itemize}
%\item Bề rộng bụng sóng là 4A, với A là biên độ của một sóng.
%\item Nếu cho biết $v_1\le v\le v_2$ thì dựa vào điều kiện sóng dừng để tìm v theo k rồi thay vào điều kiện giới hạn đó tìm ra giá trị k.
%\item Với thực nghiệm sóng dừng trên lớp ta thấy sóng dừng quan sát được với bề rộng bụng sóng lớn hơn nhiều giá trị 4A do sóng phản xạ phản xạ qua lại nhiều lần trên dây. Ở đây là chỉ xét một sóng phản xạ (bỏ qua sự phản xạ qua lại).
%\end{itemize}
%}
\end{dang}
